% sec-02: the action register and the cadence the day adopts.  Table 22,
% Figure 14 (diagrams (python)-type clustered topology).
\section{The Action Register}

Table 22 lists the day's nine actions in the order they are done. The first
reads every reply that arrived during the week, because a follow-up to an office
that has already answered would be the one mistake this day exists to prevent.
Each letter and each LinkedIn reply waits on a condition, and none is sent if its
condition fails.

\begin{apptable}
\begin{tabularx}{\textwidth}{>{\raggedright\arraybackslash}p{0.5cm} >{\raggedright\arraybackslash}p{6.4cm} >{\raggedright\arraybackslash}p{2.6cm} >{\raggedright\arraybackslash}X}
\headrow \thead{No} & \thead{Action} & \thead{File} & \thead{Waits on}\\
\midrule
1 & Read every reply of the week & The inbox & Nothing \\
2 & Write to the NCI SBIR Development Center & {\footnotesize\ttfamily email-01} & No NIH reply; NIH open \\
3 & Follow up with the SEC office & {\footnotesize\ttfamily email-02} & No reply; office open \\
4 & Follow up with the Oncology AI Program & {\footnotesize\ttfamily email-03} & No reply; office open \\
5 & Answer Applicant A & {\footnotesize\ttfamily linkedin-03} & Applicant A replied \\
6 & Answer Applicant B & {\footnotesize\ttfamily linkedin-04} & Applicant B replied \\
7 & Check the SBIR.gov registry record & {\footnotesize\ttfamily form-01} & Nothing \\
8 & Close the week's capital ledger & {\footnotesize\ttfamily capital-05} & The last session \\
9 & Adopt the weekly cadence & {\footnotesize\ttfamily brief-03} & Nothing \\
\end{tabularx}
\end{apptable}
\vspace{-0.60cm}
\tabcap{Table 22. The day's nine actions in the order they are done, each with the file that\\
carries it and the condition it waits on; no letter leaves if its condition fails.}

\subsection{The cadence the day adopts}

The last action turns the week into a rule. Figure 14 draws the weekly cadence
as a set of places and the paths a letter takes between them: in through the
inbox, the LinkedIn threads and the portals; through the 5:00 read and the 5:30
approval, where the follow-up rule is applied; out through the federal and West
Coast windows; and into the record, which the next morning reads first. Every
weekday keeps the same places and the same times, and only the theme of the two
federal windows changes from day to day \cite{autofund2026}.

\begin{appfloat}[!tb]
\begin{appfig}[diagrams (python)-type / clustered topology of one working day]
% Four clusters left to right: what comes in, the morning gate, the windows
% out, and the record.  Edges run on orthogonal buses so none crosses a label;
% one dashed edge returns from the record to the morning read.
\dgnode{dgtile}{0.00}{0.00}{\glyphmon}{Email inbox}{i1}
\dgnode{dgtile}{0.00}{-2.35}{\glyphuser}{LinkedIn threads}{i2}
\dgnode{dgtile}{0.00}{-4.70}{\glyphnet}{Portals and forms}{i3}
\dgnodew{dgtilem}{3.60}{0.00}{\glyphclock}{5:00 read, replies first}{g1}
\dgnodew{dgtilem}{3.60}{-2.35}{\glyphlock}{Follow-up rule, three open days}{g2}
\dgnodew{dgtiled}{3.60}{-4.70}{\glyphshield}{5:30 approval, one decision}{g3}
\dgnodew{dgtiled}{7.20}{0.00}{\glyphbank}{Federal window, 6:00 to 8:00}{o1}
\dgnodew{dgtiled}{7.20}{-2.35}{\glyphlink}{West Coast window, 11:00}{o2}
\dgnodew{dgtiled}{7.20}{-4.70}{\glyphbank}{Second federal window, 12:30}{o3}
\dgnodeg{dgtileg}{10.80}{0.00}{\glyphgear}{Build and commit, 2:00}{r1}
\dgnodeg{dgtileg}{10.80}{-2.35}{\glyphdoc}{Record, 3:30}{r2}
\dgnodeg{dgtileg}{10.80}{-4.70}{\glyphchart}{Friday ledger and follow-ups}{r3}
\begin{scope}[on background layer]
\node[dgcluster2,fit={(i1)(i1l)(i3)(i3l)}] (ci) {};
\node[dgcluster,fit={(g1)(g1l)(g3)(g3l)}] (cg) {};
\node[dgcluster,fit={(o1)(o1l)(o3)(o3l)}] (co) {};
\node[dgcluster2,fit={(r1)(r1l)(r3)(r3l)}] (cr) {};
\end{scope}
\node[dgctitle2,anchor=south west] at (ci.north west) {In};
\node[dgctitle,anchor=south west] at (cg.north west) {The gate};
\node[dgctitle,anchor=south west] at (co.north west) {Out};
\node[dgctitle2,anchor=south west] at (cr.north west) {The record};
\draw[dgedge] (i1) -- (g1);
\draw[dgedge,-] (i2.east) -- (1.80,-2.35) -- (1.80,0.00);
\draw[dgedge,-] (i3.east) -- (1.80,-4.70) -- (1.80,-2.35);
\draw[dgedgeb] (g1l.south) -- (g2);
\draw[dgedgeb] (g2l.south) -- (g3);
\draw[dgedgeb,-] (g3.east) -- (5.40,-4.70) -- (5.40,0.00);
\draw[dgedgeb] (5.40,0.00) -- (o1);
\draw[dgedgeb] (5.40,-2.35) -- (o2);
\draw[dgedgeb] (5.40,-4.70) -- (o3);
\draw[dgedge] (o1) -- (r1);
\draw[dgedge,-] (o2.east) -- (9.00,-2.35) -- (9.00,0.00);
\draw[dgedge,-] (o3.east) -- (9.00,-4.70) -- (9.00,-2.35);
\draw[dgedge] (r1l.south) -- (r2);
\draw[dgedge] (r2l.south) -- (r3);
\draw[dgedged] (r1.north east) -- (11.25,1.45) -- node[above,font=\tiny\sffamily] {read first, next morning} (4.05,1.45) -- (g1.north east);
\legkey{-0.60}{-6.60}{fundaccent}{The day's decisions and sends}
\legkey{6.20}{-6.60}{fundgrayl}{Inputs and the record}
\end{appfig}
\vspace{-0.60cm}
\figcaption{Figure 14. One working day as a topology of four places, with every letter passing\\
the follow-up rule and the 5:30 approval, and the record read first next morning.}
\end{appfloat}

\subsection{What the approval includes}

Approving the day approves each conditional action with its condition, not the
action alone. If the SEC office has replied, letter 2 is not sent and the reply
is answered instead. If an applicant has not replied, nothing is sent to that
applicant. If appropriations lapsed on October 1, the three federal letters wait
until their offices have been open for three business days, and nothing else on
this list changes \cite{ncisbir2026,secsf2026,fdaoceai2026}.

\actionbox{The one approval this day asks for}{Approve sending the three earned
follow-ups, each only if its condition holds; answering each applicant who has
replied, and no other; and adopting the weekly cadence, version 2, as the rule
for every week that follows.}
